% Abstract in Portuguese (resumo): single paragraph, active voice, third
% person, 150 to 500 words, followed by 3 to 5 keywords.

\begin{resumo}
  Este trabalho apresenta o desenvolvimento do Duckoder, uma plataforma \textit{web} integrada a um modelo de linguagem de grande porte (LLM) executado localmente, para auxiliar estudantes na programação de placas didáticas de sistemas embarcados. A solução é composta por uma interface \textit{web}, na qual o estudante escolhe a placa, descreve a atividade em linguagem natural e refina o código na mesma conversa; por uma API REST que insere automaticamente a descrição do \textit{hardware} da placa em cada requisição e mantém o histórico das conversas; e por um modelo aberto executado em um servidor da rede local por meio do Ollama, dispensando serviços em nuvem. A avaliação utilizou dez atividades práticas de complexidade crescente, divididas entre o ESP32 e o Arduino Mega 2560 e submetidas três vezes a dois modelos, o Qwen2.5-Coder 7B e o Llama 3.1 8B, e considerou a taxa de códigos funcionais e o tempo de resposta. O Qwen2.5-Coder produziu código funcional na primeira resposta em 83\% das gerações e em 87\% após o refinamento; o Llama 3.1, em 67\% e 77\%, respectivamente. Todas as gerações das atividades de complexidade baixa funcionaram na primeira resposta, e as falhas concentraram-se em bibliotecas específicas. O refinamento pela conversa recuperou parte dos erros de compilação, mas nenhum código que compilava e se comportava de forma diferente da pedida. O tempo médio de resposta ficou entre 12 e 14~s, e a comunicação na rede local acrescentou apenas milissegundos. Conclui-se que um modelo aberto executado na infraestrutura da instituição é suficiente para as atividades de disciplinas introdutórias, com independência de conexão externa, sem custo por requisição e com os dados mantidos na instituição.

  \vspace{\onelineskip}

  \noindent\textbf{Palavras-chave:} sistemas embarcados; modelos de linguagem de grande porte; aplicação \textit{web}; rede local; ensino.
\end{resumo}
